\documentclass[11pt,spanish]{article}
\renewcommand{\sfdefault}{lmss}
\renewcommand{\familydefault}{\rmdefault}
\usepackage[T1]{fontenc}
\usepackage[utf8]{inputenc}
\usepackage{babel}
\addto\shorthandsspanish{\spanishdeactivate{~<>.}}
\deactivatequoting
\usepackage{geometry}
\geometry{tmargin=2cm,bmargin=2cm,lmargin=2cm,rmargin=2cm}
\usepackage{fancyhdr}
\pagestyle{fancy}
\setlength{\headheight}{14pt}
\usepackage[skip=0.55\baselineskip]{parskip}
\usepackage{microtype}
\usepackage{hyperref}
\hypersetup{hidelinks}
\usepackage{xcolor}
\usepackage{enumitem}
\usepackage{listings}
\usepackage{needspace}
\usepackage{titling}
\setlength{\droptitle}{-2cm}

\definecolor{rojo7a3c3d}{rgb}{0.478,0.235,0.239}
\definecolor{verde587246}{rgb}{0.345,0.447,0.275}
\definecolor{codegray}{rgb}{0.96,0.96,0.96}

\newcommand{\cf}[2]{\colorbox{#1}{\textcolor{white}{\strut #2}}}
\newcounter{ejercicio}
\newcommand{\ejercicio}[1]{%
  \stepcounter{ejercicio}%
  \textbf{Ejercicio \theejercicio. #1}%
}

\lstset{
  basicstyle=\ttfamily\small,
  backgroundcolor=\color{codegray},
  frame=single,
  rulecolor=\color{black!25},
  breaklines=true,
  showstringspaces=false,
  columns=fullflexible,
  aboveskip=0.35em,
  belowskip=0.25em,
  literate={á}{{\'a}}1 {é}{{\'e}}1 {í}{{\'i}}1 {ó}{{\'o}}1 {ú}{{\'u}}1 {ñ}{{\~n}}1
}

\setlength{\parindent}{0pt}
\setlist[enumerate]{leftmargin=*,itemsep=0.25em,topsep=0.3em}
\setlist[itemize]{leftmargin=1.4em,itemsep=0.15em,topsep=0.25em,parsep=0pt}

\begin{document}

\lhead{PDS}
\rhead{Práctica 01}
\title{Programación de Sistemas\\ \bigskip
\cf{rojo7a3c3d}{\textbf{\Large Práctica 01: Extensión del analizador léxico para \textsc{IMP++}}}}
\author{Manuel Soto Romero \and Araceli Liliana Reyes Cabello}
\date{Semestre 2026-2 \\ Colegio de Ciencia y Tecnología UACM}

\maketitle

\section*{\cf{verde587246}{Descripción}}

Extiendan el analizador léxico funcional de \textsc{IMP} desarrollado en clase
para reconocer las construcciones de \textsc{IMP++}. Deberán modificar sus
reglas léxicas, conservar la compatibilidad con \textsc{IMP} y comprobar el
resultado con pruebas. Esta versión será la entrada de la Práctica 2.

La práctica se realiza en equipos de una a tres personas.
\ifdefined\incluirepositorio
Cada equipo entregará su trabajo mediante un repositorio privado compartido.
\fi

\section*{\cf{verde587246}{Objetivos}}

\begin{enumerate}
  \item Extender la especificación léxica de \textsc{IMP} con los tokens de \textsc{IMP++}.
  \item Resolver conflictos entre palabras reservadas, identificadores y operadores compuestos.
  \item Verificar la extensión sin romper los programas válidos de \textsc{IMP}.
\end{enumerate}

\section*{\cf{verde587246}{Desarrollo}}

\ifdefined\incluirepositorio
\ejercicio{Creen el repositorio del equipo.}
\begin{itemize}
  \item Formen un equipo de una a tres personas y acuerden un nombre para identificarlo.
  \item Una persona del equipo debe crear en \textsc{GitHub} un repositorio privado para las cuatro prácticas.
  \item Añadan como colaboradores a todas las personas del equipo y al usuario \texttt{manu-msr}.
  \item Registren en el \texttt{README.md} el nombre del equipo, los nombres de sus integrantes y sus usuarios de \textsc{GitHub}.
  \item Copien la carpeta \path{scripts_estudiante} al repositorio y renómbrenla como \path{practica1_lexer_imp_mas_mas}.
  \item Al terminar, el acceso registrado en \textsc{GitHub} debe cumplir:
\begin{lstlisting}[language={}]
Visibilidad: privada
Equipo: 1 a 3 integrantes con acceso
manu-msr: invitación pendiente o colaborador
\end{lstlisting}
  \item Para entregar, una persona del equipo debe pegar la liga al repositorio en los comentarios de la actividad en \textsc{Google Classroom}.
\end{itemize}
\fi
\Needspace{6\baselineskip}
\ejercicio{Verifiquen el punto de partida.}
\begin{itemize}
  \item Preparen el entorno con las instrucciones del \texttt{README.md}; todavía no modifiquen el lexer.
  \Needspace{4\baselineskip}
  \item Ejecuten el ejemplo base y las pruebas:
\begin{lstlisting}[language={}]
python lexer.py ejemplos/imp_base.imp
python -m unittest -v
\end{lstlisting}
  \Needspace{10\baselineskip}
  \item Deben aparecer 14 tokens. Ejemplo de salida:
\begin{lstlisting}[language={}]
LOC        'X'              1:1
ASSIGN     ':='             1:3
NUM        '0'              1:6
SEMI       ';'              1:7
...
WHILE      'while'          2:1
...
Ran 10 tests in ...
OK (skipped=5)
\end{lstlisting}
  \item Si el inicio, la posición de \texttt{while} o el resumen no coinciden, revisen la instalación o la copia de los archivos.
\end{itemize}

\Needspace{6\baselineskip}
\ejercicio{Extiendan identificadores y comentarios.}
\begin{itemize}
  \item Abran \texttt{imp.lark} y modifiquen \texttt{LOC} con el patrón \texttt{[A-Za-z][A-Za-z0-9\_]\allowbreak*}. Comprueben que \texttt{contador\_1} se reconozca como un solo \texttt{LOC}.
  \item Definan \texttt{COMMENT} para reconocer desde \texttt{//} hasta el final de la línea y agreguen la instrucción~\texttt{\%ignore COMMENT}.
  \item Retiren los decoradores \texttt{@unittest.skip} de las pruebas de identificadores y comentarios.
  \Needspace{4\baselineskip}
  \item Ejecuten las pruebas:
\begin{lstlisting}[language={}]
python -m unittest -v
\end{lstlisting}
  \Needspace{7\baselineskip}
  \item Ejemplo de salida:
\begin{lstlisting}[language={}]
test_comentario_se_ignora ... ok
test_identificador_con_guion_bajo ... ok
...
Ran 10 tests in ...
OK (skipped=3)
\end{lstlisting}
  \item Estas pruebas revisan que \texttt{contador\_1} produzca un solo \texttt{LOC}, que el comentario no produzca tokens y que \texttt{Y} conserve la posición línea 2, columna 1.
\end{itemize}

\Needspace{6\baselineskip}
\ejercicio{Agreguen palabras y símbolos de \textsc{IMP++}.}
\begin{itemize}
  \item Agreguen en \texttt{imp.lark} las reglas para \texttt{int}, \texttt{bool}, \texttt{print} y \texttt{for}. Conserven los nombres de token indicados en el archivo.
  \item Usen la misma prioridad y el mismo límite léxico de las palabras reservadas de \textsc{IMP}.
  \item Agreguen las reglas para las llaves y para los operadores \texttt{++} y \texttt{-{}-}.
  \Needspace{4\baselineskip}
  \item Retiren los decoradores \texttt{@unittest.skip} restantes y ejecuten las pruebas:
\begin{lstlisting}[language={}]
python -m unittest -v
\end{lstlisting}
  \Needspace{8\baselineskip}
  \item Ejemplo de salida:
\begin{lstlisting}[language={}]
test_incremento_y_decremento_son_tokens_completos ... ok
test_palabras_reservadas_y_llaves ... ok
test_programa_imp_mas_mas ... ok
...
Ran 10 tests in ...
OK
\end{lstlisting}
  \item Las pruebas comparan, entre otros casos, \texttt{print printable} con \texttt{PRINT, LOC} y \texttt{X++ + Y-{}- - 1} con \texttt{LOC, INC, PLUS, LOC, DEC, MINUS, NUM}.
\end{itemize}

\Needspace{6\baselineskip}
\ejercicio{Comprueben la compatibilidad y documenten.}
\begin{itemize}
  \Needspace{6\baselineskip}
  \item Ejecuten el ejemplo completo y todas las pruebas:
\begin{lstlisting}[language={}]
python lexer.py ejemplos/imp_mas_mas.imp
python lexer.py ejemplos/error_lexico.imp
python -m unittest -v
\end{lstlisting}
  \Needspace{9\baselineskip}
  \item \texttt{imp\_mas\_mas.imp} debe producir 30 tokens y no debe mostrar \texttt{COMMENT}. Ejemplo de salida:
\begin{lstlisting}[language={}]
INT        'int'            1:1
LOC        'contador_1'     1:5
ASSIGN     ':='             1:16
NUM        '0'              1:19
SEMI       ';'              1:20
...
PRINT      'print'          5:5
LPAR       '('              5:10
LOC        'contador_1'     5:11
RPAR       ')'              5:21
SEMI       ';'              5:22
RBRACE     '}'              6:1
\end{lstlisting}
  \item Ejemplo de salida para el archivo con error:
\begin{lstlisting}[language={}]
Error léxico en línea 2, columna 8: '@'
\end{lstlisting}
  \item Agreguen una prueba que combine al menos dos elementos nuevos.
  \item Agreguen otra para un carácter no reconocido y comprueben que produzca \texttt{UnexpectedCharacters}.
  \item Ejecuten nuevamente las pruebas. Ejemplo del resumen después de agregar sólo los dos casos solicitados:
\begin{lstlisting}[language={}]
...
Ran 12 tests in ...
OK
\end{lstlisting}
  \item Completen el \texttt{README.md} con el resultado, una decisión léxica y el tiempo real de trabajo.
  \item Entreguen la carpeta sin pruebas omitidas ni tareas \texttt{TODO} pendientes.
\end{itemize}

\section*{\cf{verde587246}{Referencias de consulta}}

\begin{enumerate}[label=\textbf{[\arabic*]}]
  \item Winskel, G. (1993). \textit{The Formal Semantics of Programming Languages: An Introduction}, sección 2.1. The MIT Press.
  \item Lark. \textit{Grammar Reference y API Reference}. \url{https://lark-parser.readthedocs.io/}
\end{enumerate}

\end{document}
